For each of the following systems, find out how many solutions they have. 
Check also if there are redundant equations that can be removed.

\paragraph{Hint} In order to identify the rank of the matrix, you may want to
transform it into a row echelon form using elementary row operations. A matrix is in row echelon form  when it satisfies the following conditions:
\begin{itemize}
\item each leading entry is in a column to the right of the leading entry in the previous row,
\item rows with all zero elements, if any, are below rows having a non-zero element.
\end{itemize}
Some definitions also impose that the first non-zero element in each
row is 1. However, this conditions is not needed to obtain the rank of
the matrix. 

\subsubsection*{Question 1}
\begin{align*}
x_1+2x_2+3x_3 & = 0,\\
4x_1+5x_2+6x_3 & = 0,\\
7x_1+8x_2& =0.
\end{align*}
\subsubsection*{Question 2}
\begin{align*}
x_1+2x_2+3x_3 & = 0,\\
4x_1+5x_2+6x_3 & = 0,\\
7x_1+8x_2+9x_3 & =0.
\end{align*}
\subsubsection*{Question 3}

\begin{align*}
x_1+2x_2+3x_3 & = 0,\\
2x_1+4x_2+6x_3 & =0,\\
5x_1+10x_2+15x_3 & =0.
\end{align*}
\subsubsection*{Question 4}

\begin{align*}
x_1 + x_2 + x_3 &= 1,\\
x_1 - x_2 + x_4 &= 1, \\
x_1 -5x_2 -2x_3 +3x_4 &= 1.
\end{align*}
\subsubsection*{Question 5}

\begin{align*}
x_1+x_2&=0,\\
x_1+2x_2+2x_3&=0,\\
3x_2+3x_3&=0.
\end{align*}
